\section{Introduction}\label{sec:introduction65}
A learned continuation is economically useful when it changes what can be computed and which decisions can be justified under the original law. Neural Bellman Operators (NBOs) connect a constructed continuation, its feasible action witnesses, and the complete policy actually implemented. The central comparison is
\[
 \underline V\leq V^*\leq J^{\widehat\pi}\leq\overline J.
\]
The lower inequality covers every original feasible action. The upper inequality includes observation, numerical action realization, and the controller's subsequent choices. A guarantee for a fitted surrogate is not a substitute for this economic comparison.

Our first contribution is constructive. Each dated NBO is formed against its own already attained future, using primitive state and action covers, integration allowances, and feasible action owners. The affine--ReLU realization does not assume successful unrestricted nonconvex training. Original-owner compilation preserves the continuation and its feasible witness while changing evaluation work. Joint-domain and common-readout extensions retain the coupling between decisions at different acquired states. Their hypotheses and complete statements remain in the accompanying material.

Our second contribution is an original-law account for complete policies under recursive risk. Centering the original conditional dual at the stored continuation retains the penalty for changing the maximizing risk density. State-dependent subsolutions and supersolutions then compare the implemented policy with the original optimum. The continuous investment construction derives curvature, value supports, and all-action inequalities from the original technology. Feasible risk witnesses minorize the lower action problem; full-law envelopes dominate the policy value. Representation quality, certificate validity, and computational gain therefore have distinct, testable meanings.

Our third contribution joins information resolution to neural acquisition. For convex Lipschitz continuations, the continuous-law AVaR bracket has second-order innovation resolution. We establish a matching exponent against all deterministic adaptive function-evaluation procedures, rather than only sharpness of the displayed quadrature pair. An exact oblique-ridge comparison isolates a representation advantage over additive approximations and quadratics; it does not imply an advantage over a conventional method supplied with the same hinge geometry. A certificate-oriented readout program gives an implementable construction whose retained whole-domain residual span cannot be worse than that of its least-squares control. Numerical optimizer success is not used as a certificate.

The empirical design distinguishes these mathematical mechanisms. The retained R96 comparison fixes six-dimensional design rank, 1,024 labels, and a 128-byte payload before 4,096 evaluation states at each of two risk levels. Its adverse result remains central: quadratic storage is faster than ReLU storage, and equivalence at the specified margin is not an economic improvement. A new fixed 84-configuration campaign varies rank and label budgets and compares least-squares and certificate-oriented readouts within each representation. It records the first proposal, rejection reasons, paid repair, refinements, exact fallback, and all acquisition stages. These new construction diagnostics are not relabeled as a second population confirmation of services already selected using the earlier evidence.

The statistical estimand is the mean under a uniform 20-bit dyadic initial-state distribution; the construction domain and subsequent innovation law are continuous. Qualification of the complete service supplies the support of paired differences. The R96 engineering margin is one percent of the maximal first-date investment charge. That normalization is not an empirical welfare calibration. Repeated clocks add no independent policy observations. Independent exact arithmetic checks the retained paired inference, and a separately derived implementation replays conventional as well as fitted proposals; the scope of exact endpoint checks is reported separately.

Finally, an economic purchase must pay for acquisition and the \emph{whole} deployed workload. We give a unit-consistent interval frontier that multiplies the per-decision policy contrast by the number of decisions in a timed workload. The measured acquisition bracket includes primitive labels, search, fitting, domain validation, serialization, checking, and a conservative full-process upper charge. Computing a difference of such brackets requires interval subtraction, not subtraction of two upper endpoints. Workload clocks supply conditional sensitivity inputs, not confidence bounds for future machines. Externally supplied prices and fixed fees remain explicit arguments of the purchase decision.

\paragraph{Related work.}
Lipschitz extension and distance transforms are established construction tools \citep{mcshane1934,felzenszwalb2012}. Approximate policy iteration, safe improvement, and convex dynamic programming have established comparison results \citep{munos2003,pirotta2013,scherrer2015,yang2020convex}. The NBO construction retains the attained own-future approximation and feasible owner throughout the primitive accuracy account. A native and a neural encoding of the same function receive the same economic comparison.

Risk-averse Bellman equations and feature-based risk learning are established directions \citep{ruszczynski2010risk,ruszczynski2014erratum,ruszczynskiyang2025}. AVaR threshold optimization, convex quadrature inequalities, and Hoeffding concentration are not new here \citep{RockafellarUryasev2000,alomari2016hermite,hoeffding1963}. The information lower bound uses the classical indistinguishable-input principle \citep{Novak2015Integration97}, with an explicit convex tail construction suited to the original-law policy problem. The comparison table in Section~\ref{sec:comparison95} locates the retained qualification results relative to their mathematical ingredients.

\paragraph{Organization.}
The article proceeds from own-future construction to complete-policy and recursive-risk qualification, continuous investment, information and representation, and the two distinct empirical catalogues. The focused Technical Supplement contains proofs used by the current article. The Numerical and Development Appendix retains the primitive fixed-accuracy investment account, horizon allocation, intervening proofs, broader applications, and adverse vector-scaling evidence as complete sections. This division preserves the research record while distinguishing current results from historical support; it does not change the NBO problem.
